% ======================================================================
% sections/results.tex
%
% Results scaffold for the paper:
%   Triple Humanoid 24/7 Adverse Event Oncology Trial Response Team:
%   4-Site Rotation
%
% This file is a draft scaffold. A future Claude Code Opus 4.7 1M Max
% session reads the bracketed instruction blocks below, then reads the
% files those instructions name, then expands the scaffold into the
% finished results section. Target length is roughly 14 to 18 pages of
% the 70 plus page final paper.
% ======================================================================

\begin{sectionaccent}

[RESULTS INSTRUCTIONS FOR FUTURE CLAUDE CODE OPUS 4.7 1M MAX SESSION.
Write the results section as 7 subsections inside this single
\texttt{\textbackslash section\{Results\}} block. Every numeric claim must
come from a file named below. Do not introduce numbers that do not
appear in those files. Quote file paths in the prose so a reader can
verify each number.

Subsection 1, Per-Site LLM Decision Streams. Read the 4 site files in
\path{demo-projects/07-humanoid/paper/execution/data/site_runtime_runs/}:
\begin{itemize}
  \item \path{AT-01_llm_decisions.jsonl} (Atlanta, 62,330 bytes).
  \item \path{BO-01_llm_decisions.jsonl} (Boston, 62,330 bytes).
  \item \path{SD-01_llm_decisions.jsonl} (San Diego, 62,330 bytes).
  \item \path{SF-01_llm_decisions.jsonl} (San Francisco,
        62,330 bytes).
\end{itemize}
All 4 site files have the same byte length. Treat the byte length
match as a parallel-scope check: the 1 Hz broadcast cadence and the
shared schema produce decision logs of the same record count per
site. Read every line of \path{AT-01_llm_decisions.jsonl} (it is
the lead Atlanta file flagged as Hz data at scale for the project).
Then spot-check the other 3 site files for record count parity and
schema parity. If the records repeat across sites, call that out as a
duplication artifact in the runtime that production should fix
without retracting the parallel-scope claim. If the records differ
across sites despite identical byte length, call that out as
acceptable schema parity. Do this check explicitly and report the
finding in the prose.

Subsection 2, Iteration Sweep Aggregate. Read:
\begin{itemize}
  \item \path{demo-projects/07-humanoid/paper/codegen/data/iterations/index.jsonl}
        for the 32 iteration codegen index. State the 32 iteration
        count and the 4 site coverage.
  \item \path{demo-projects/07-humanoid/paper/execution/data/iterations_index_original.jsonl}
        and
        \path{demo-projects/07-humanoid/paper/execution/data/iterations_index_runtime.jsonl}
        for the original and runtime iteration indexes.
  \item \path{demo-projects/07-humanoid/paper/execution/data/iterations_aggregate.duckdb}
        as the DuckDB roll up. Discuss that DuckDB and JSONL coexist
        cleanly in the execution data layer. The DuckDB plus JSONL
        pairing demonstrates DuckDB proficiency for analytical queries
        and JSONL proficiency for append-only logs.
\end{itemize}

Subsection 3, Comparison Reports. Read:
\begin{itemize}
  \item \path{demo-projects/07-humanoid/paper/codegen/reports/comparison.json}.
  \item \path{demo-projects/07-humanoid/paper/codegen/reports/comparison_methodology.md}.
  \item \path{demo-projects/07-humanoid/paper/codegen/reports/report.md}.
  \item \path{demo-projects/07-humanoid/paper/codegen/reports/dashboard.html}.
        Treat the dashboard as the primary interpretability artifact
        for future paper viewers; the rendered HTML adds project
        interpretability without requiring the reader to inspect raw
        JSON. Embed a screenshot reference in the prose and link to
        the file path.
  \item \path{demo-projects/07-humanoid/paper/execution/reports/comparison_v0_4_0_copy.json}.
  \item \path{demo-projects/07-humanoid/paper/execution/reports/comparison_methodology_v0_4_0_copy.md}.
  \item \path{demo-projects/07-humanoid/paper/execution/reports/dashboard_v0_4_0_copy.html}.
  \item \path{demo-projects/07-humanoid/paper/execution/reports/execution_report_v0_5_0.md}.
  \item \path{demo-projects/07-humanoid/paper/execution/reports/report_v0_4_0_copy.md}.
\end{itemize}

Subsection 4, Comparison ASCII Visuals. Include verbatim inside
\texttt{asciibox} environments:
\begin{itemize}
  \item \path{demo-projects/07-humanoid/paper/execution/diagrams/04_comparison_radar.txt}
        as Figure 7.
  \item \path{demo-projects/07-humanoid/paper/execution/diagrams/01_module_dependency.txt}
        as Figure 8 if it was not already placed in the Methods
        section. If it was already placed there, only cross reference.
\end{itemize}

Subsection 5, Output Markdown Tables. Read:
\begin{itemize}
  \item \path{demo-projects/07-humanoid/paper/execution/outputs/comparison_summary_table.md}.
  \item \path{demo-projects/07-humanoid/paper/execution/outputs/iteration_sweep_table.md}.
  \item \path{demo-projects/07-humanoid/paper/execution/outputs/execution_tree.txt}.
  \item \path{demo-projects/07-humanoid/paper/execution/outputs/file_inventory.txt}.
\end{itemize}
Translate each markdown table into a \texttt{\textbackslash sectiontable}
LaTeX table using the 8 column shared macro. Do not change the data;
only re-render the format.

Subsection 6, Figures. Reference every figure in
\path{demo-projects/07-humanoid/paper/execution/figures/} (6 PNGs at
300 dpi: swarm architecture, response time histogram, camaraderie
heatmap, role rotation timeline, force budget, 4 site comparison).
Use \texttt{\textbackslash includegraphics} only if the figure file is
copied into the draft-paper folder as part of the future expansion.
Otherwise reference the file path in a caption note.

Subsection 7, Headline Metrics. Pull the headline numbers into one
final \texttt{\textbackslash sectiontable}:
\begin{itemize}
  \item Sites: 4 (SF-01, SD-01, BO-01, AT-01).
  \item Robots per site: 3 Unitree H2 EDU.
  \item Robots total: 12.
  \item LLM cadence: 1 Hz broadcast.
  \item Motion cadence: 10 Hz per robot.
  \item Monitoring window: 168 hours.
  \item AE volume: about 84 across the week.
  \item SAE volume: about 24 (CTCAE grade 3 or higher).
  \item Iterations: 32 deterministic sweeps.
  \item E-stop latency: 5 ms swarm wide.
  \item Per-arm cumulative force: 15 N during chest compressions,
        5 N otherwise.
  \item Inter-robot minimum distance: 1.2 m at rest, 0.4 m during
        shared task handoff.
  \item Cumulative cross-robot force: 22 N during 3-robot patient
        transfer.
  \item Camaraderie reduction factor: 3 in single robot error
        potential.
\end{itemize}
Every number in this final table must trace to a file in
\path{instructions/}, \texttt{codegen/}, or \texttt{execution/}.

Required real-life feasibility insight to include in this section:
``Identical byte length across the 4 site decision logs is not the
production end state. In a real 168 hour run the per-site record
counts will diverge because each site sees a different AE arrival
distribution. The current parity demonstrates that the runtime
stamps every site at 1 Hz on a shared schema; the next iteration is
to feed each site a site-specific arrival stream.''

Do not exceed 18 pages. Use single dashes only. Use \S\ for the
section sign. Avoid stranding 1 or 2 word lines on their own.]

The results section will present the per-site LLM decision streams,
the 32 iteration sweep aggregate, the comparison reports and ASCII
visuals, the markdown output tables converted to LaTeX tables, the
figures, and the final headline metrics table.

\sectiontable{Results planning matrix}{%
Decisions & R1 & 1 & runtime runs & 4 site walk & 2 pages & 62,330 bytes & Ready \\
Iterations & R2 & 2 & index, duckdb & 32 sweep & 2 pages & DuckDB JSONL & Ready \\
Reports & R3 & 3 & code, exec reports & Comparison & 3 pages & dashboard named & Ready \\
ASCII & R4 & 4 & exec diagrams & Figures 7, 8 & 1 to 2 pages & Verbatim text & Ready \\
Tables & R5 & 5 & exec outputs & Md to LaTeX & 2 to 3 pages & 8 col macro & Ready \\
Figures & R6 & 6 & exec figures & 6 PNG callouts & 1 to 2 pages & 300 dpi & Ready \\
Metrics & R7 & 7 & instructions metrics & Headline table & 1 page & Every number sourced & Ready \\
}

\end{sectionaccent}
